% 两种世界的可见源样本和目标输入相同；目标标签不作为算法输入。
% 两域输入均匀分布在{0,1}，源标签y=x；目标标签可能为x或1-x。
% 固定同一算法随机数后输出相同；此处示例h(0)=0,h(1)=1。
\begin{tikzpicture}[foundation picture]
 \path[use as bounding box] (0,0) rectangle (112,77);
 \node at (19,73) {可见数据};
 \draw[foundation frame,fill=white] (0,35) rectangle (38,68);
 \node at (19,62) {源样本$S_S$};
 \node at (19,55) {$(0,0),\ (1,1)$};
 \node at (19,46) {目标输入$U_T$};
 \node at (19,39) {$(0,?),\ (1,?)$};
 \node[foundation model,fill=white,minimum width=38mm,minimum height=9mm] (algorithm)
  at (19,23) {$h=A(S_S,U_T)$};
 \draw[foundation flow] (19,35)--(algorithm.north);
 \node[foundation output,fill=white,minimum width=38mm,minimum height=9mm] (prediction)
  at (19,8) {$h(0)=0,\quad h(1)=1$};
 \draw[foundation flow] (algorithm.south)--(prediction.north);
 % 唯一输出复制到两个真值比较；隐藏标签没有返回算法的边。
 \draw[FigureStroke,line width=1.1pt] (prediction.east)--(47,8)--(47,54.5);
 \draw[foundation flow] (47,54.5)--(54,54.5);
 \draw[foundation flow] (47,17.5)--(54,17.5);
 \foreach \base/\case/\yzero/\yone in {43/1/0/1,6/2/1/0}{
  \node at (83,\base+28) {隐藏目标标签\case};
  \draw[foundation frame,fill=white] (54,\base) rectangle (112,\base+23);
  \node at (70,\base+18.5) {$x=0$};
  \node at (97,\base+18.5) {$x=1$};
  \node at (70,\base+11.5) {$h(0)=0$};
  \node at (97,\base+11.5) {$h(1)=1$};
  \node at (70,\base+4.5) {$y=\yzero$};
  \node at (97,\base+4.5) {$y=\yone$};
 }
 \foreach \x in {81,108}{
  \node[text=FigureBlue!65!black] at (\x,47.5) {$\checkmark$};
  \node[text=FigureRed!70!black] at (\x,10.5) {$\times$};
 }
\end{tikzpicture}%
